% ===========================================================================
% track 收尾 (1/2):Track Lifecycle 與 Birth
% ===========================================================================
\chapter{Track Lifecycle 與 Birth}\label{ch:lifecycle}

\begin{center}\archmap{track}\end{center}

\begin{contract}{Lifecycle(track 子機制)}
\textbf{輸入}\; assignment 結果(matched / unmatched track / unmatched detection)\quad
\textbf{輸出}\; 更新後的 track 集合與狀態(tentative/confirmed/lost/removed)\\
\textbf{方法}\; hit-streak 累積 confirm、age 計數 lost/remove\quad
\textbf{Source}\; \texttt{tracker\_gpu.cu}\\
\textbf{Baseline}\; \texttt{new\_track\_thresh: 0.28}、\texttt{confirm\_streak: 3}、
\texttt{confirm\_score\_thresh: 0.50}、\texttt{track\_buffer: 30}、\texttt{per\_seq\_adapt: false}
\end{contract}

\paragraph{演算法傳遞.} assignment 後依配對狀態驅動 track 狀態機;下圖為主要轉移
(tentative 若連續未配會很快 drop,圖中略):

\begin{center}\resizebox{\linewidth}{!}{%
\begin{tikzpicture}[node distance=20mm]
  \node[step=Ctrk] (nw) {未配 det};
  \node[step=Ctrk, right=of nw] (te) {tentative};
  \node[step=Ctrk, right=of te] (co) {confirmed};
  \node[step=Ctrk, right=of co] (lo) {lost};
  \node[step=Ctrk, right=of lo] (rm) {removed};
  \draw[flowarr=Ctrk] (nw) -- node[flowlbl,above]{$s_j{\ge}$new\_thr} (te);
  \draw[flowarr=Ctrk] (te) -- node[flowlbl,above]{streak${\ge}3$}    (co);
  \draw[flowarr=Ctrk] (co) -- node[flowlbl,above]{unmatched}         (lo);
  \draw[flowarr=Ctrk] (lo) -- node[flowlbl,above]{age${>}$buffer}    (rm);
  \draw[flowarr=Ctrk] (lo) to[bend left=32] node[flowlbl,above]{re-match} (co);
\end{tikzpicture}}\end{center}

\section{State 轉移}\label{sec:lc-states}
assignment 後 tracker 更新 state:
\begin{itemize}\setlength\itemsep{2pt}
  \item \textbf{matched}(confirmed/tentative):以 detection \(z_j\) 做 Kalman update
        (\Cref{ch:kalman}),\texttt{hit\_streak}\(+1\),\texttt{age}/\texttt{time\_since\_update}
        歸零。
  \item \textbf{unmatched active track}:\texttt{age}\(\le\)\texttt{track\_buffer} 期間維持
        active/lost,超過 max age 後 remove。
  \item \textbf{unmatched detection}(\(s_j\ge\)\texttt{new\_track\_thresh}):生成新的
        tentative track。
\end{itemize}

\section{Birth 與 Confirm}\label{sec:lc-confirm}
新 tentative track 必須累積足夠連續 evidence 才會被視為 confirmed output:
\begin{equation}\label{eq:lc-confirm}
\begin{aligned}
  \text{confirmed}\iff{}&
  \texttt{hit\_streak}\ge\texttt{confirm\_streak}\\
  &{}\land\bar{s}\ge\texttt{confirm\_score\_thresh}.
\end{aligned}
\end{equation}
baseline:
\begin{center}\small
\begin{tabular}{@{}ll@{\qquad}ll@{}}
\texttt{confirm\_streak} & 3 &
\texttt{confirm\_score\_thresh} & 0.50 \\
\texttt{new\_track\_thresh} & 0.28 &
\texttt{track\_buffer} & 30
\end{tabular}
\end{center}
若干 birth-gate experiments 位於 \path{scripts/eval/config/lifecycle.py},但目前
preset 未啟用(\texttt{per\_seq\_adapt: false})。
普通新生 track 以 tentative 與 \texttt{hit\_streak=1} 起跑;bridge relink 復活的
slot 會保留 lost id 並直接回到 confirmed,不再走 tentative confirm gate。

\suppinput{supp/09-lifecycle-impl}
